\chapter{Beyond NP}\label{ch:beyondnp}

\begin{center}\itshape
Guessing one solution is the first rung of a ladder;\\
playing a game well, or counting all the solutions, lies far above.\\
\upshape(the content of this chapter, in our words)
\end{center}

A puzzle asks: is there a way? A game asks something harder: is there a move, such that whatever my opponent answers, there is a move, such that whatever they answer\dots\ and so on until the end. In 1972 and 1973, at MIT, Albert Meyer and his student Larry Stockmeyer studied what happens when the quantifiers of a problem alternate. Studying the problem of deciding whether two regular expressions describe the same language, they defined a whole hierarchy of classes above $\mathsf{NP}$, one for each number of alternations, and found that formulas with an unbounded number of alternating quantifiers are complete for polynomial \emph{space}. Games joined at once: Thomas Schaefer showed in 1978 that a children's word game, in which players name places whose first letter is the last letter of the previous one, becomes $\mathsf{PSPACE}$-complete when played on an arbitrary graph. Go on an $n\times n$ board was shown $\mathsf{PSPACE}$-hard by David Lichtenstein and Michael Sipser in 1980, and perfect play at chess on an $n\times n$ board needs exponential time, by a theorem of Aviezri Fraenkel and Lichtenstein of 1981.

Counting is another way up. In 1913 George Pólya asked whether one could change the signs of some entries of a matrix so that its determinant becomes its \emph{permanent}, the same sum over permutations without the signs. The determinant is computed quickly by Gaussian elimination; the permanent of a $0$--$1$ matrix counts the perfect matchings of a bipartite graph. In 1979 Leslie Valiant explained why nobody had found a fast algorithm for it: computing the permanent is complete for the class $\#\mathsf P$ of counting problems, although deciding whether a perfect matching exists is easy. Twelve years later Seinosuke Toda proved that counting is so powerful that, with a counting oracle, one can decide everything in Meyer and Stockmeyer's hierarchy. Valiant received the Turing Award in 2010, and Toda the Gödel Prize in 1998.

The chapter ends with a modern turn. NP-hardness says that an algorithm is probably not polynomial; it does not say whether quadratic time can be improved to $n^{1.9}$. Since 2001 a precise conjecture about satisfiability, the strong exponential time hypothesis, has been used to explain why decades of effort did not improve some quadratic algorithms, such as the edit distance algorithm of Tome~C2, Chapter~5.

\begin{think}
In a game of noughts and crosses on an $n\times n$ board, where $k$ in a row wins, what would you need to remember to decide who wins with perfect play? How long would it take?
\end{think}

\section{The polynomial hierarchy}\label{sec:ph}

By Proposition~\ref{prop:certificates}, a language is in $\mathsf{NP}$ when membership is expressed by one existential quantifier over short certificates, followed by a polynomial-time check; in $\mathsf{coNP}$ when it is expressed by one universal quantifier. Allowing more quantifiers gives a hierarchy.

\begin{definition}\label{def:ph}
For $k\ge1$, a language $L$ is in $\Sigma_k^p$ if there are a polynomial $p$ and a polynomial-time machine $V$ such that
\[
x\in L\iff\exists y_1\,\forall y_2\,\exists y_3\cdots Q_ky_k\ \ V\text{ accepts }\langle x,y_1,\dots,y_k\rangle,
\]
where the quantifiers alternate, $Q_k$ is $\exists$ or $\forall$ according to the parity of $k$, and each $y_i$ ranges over $\{0,1\}^{p(|x|)}$. The class $\Pi_k^p$ is defined in the same way with the quantifiers starting with $\forall$, so that $\Pi_k^p=\mathsf{co}\Sigma_k^p$. We set $\Sigma_0^p=\Pi_0^p=\mathsf P$ and $\mathsf{PH}=\bigcup_k\Sigma_k^p$, the \emph{polynomial hierarchy}.
\end{definition}

So $\Sigma_1^p=\mathsf{NP}$ and $\Pi_1^p=\mathsf{coNP}$. A typical problem of $\Sigma_2^p$ is the minimisation of circuits: given a Boolean formula $\varphi$ and a number $k$, is there a formula of size at most $k$ equivalent to $\varphi$? It reads ``there is a formula $\psi$ such that for every assignment $a$, $\varphi(a)=\psi(a)$'' (Exercise~\ref{ex:minformula}). Christopher Umans proved in 1998 that a version of it, for formulas in disjunctive normal form, is complete for $\Sigma_2^p$ (quoted).

\begin{proposition}\label{prop:phbasic}
For every $k$, $\Sigma_k^p\cup\Pi_k^p\subseteq\Sigma_{k+1}^p\cap\Pi_{k+1}^p$, and $\mathsf{PH}\subseteq\mathsf{PSPACE}$.
\end{proposition}

\begin{proof}
A quantifier over a variable that $V$ ignores can be added anywhere: in front of the formula of $\Sigma_k^p$ with the opposite quantifier, or after it with the right parity. This gives the inclusions. For $\mathsf{PSPACE}$: to decide $x\in L$ for $L\in\Sigma_k^p$, loop over $y_1$, inside it over $y_2$, and so on, evaluating the quantifiers recursively; this takes space $O(kp(|x|))$ for the $y_i$, plus the space of $V$, reused for every tuple.
\end{proof}

Figure~\ref{fig:hierarchy} shows the hierarchy. All its inclusions are believed to be strict, but none is known to be: for all we know, $\mathsf P=\mathsf{PSPACE}$. What is known is that the hierarchy cannot be strict for a while and then stop growing by accident: equality at one level propagates upwards.

\begin{figure}[ht]
\centering
\fig{ch10-hierarchy}
\caption{The polynomial hierarchy: each level adds one alternation of quantifiers. Every arrow is an inclusion; none is known to be strict.}
\label{fig:hierarchy}
\end{figure}

\begin{theorem}\label{thm:collapse}
If $\Sigma_k^p=\Pi_k^p$ for some $k\ge1$, then $\mathsf{PH}=\Sigma_k^p$. In particular, if $\mathsf P=\mathsf{NP}$ then $\mathsf{PH}=\mathsf P$.
\end{theorem}

\begin{proof}
We show $\Sigma_{k+1}^p\subseteq\Sigma_k^p$; then $\Pi_{k+1}^p=\mathsf{co}\Sigma_{k+1}^p\subseteq\Pi_k^p=\Sigma_k^p$, and by induction every level is $\Sigma_k^p$. Let $L\in\Sigma_{k+1}^p$:
\[
x\in L\iff\exists y_1\ \langle x,y_1\rangle\in L',
\]
where
\[
\langle x,y_1\rangle\in L'\iff\forall y_2\,\exists y_3\cdots V\text{ accepts }\langle x,y_1,\dots,y_{k+1}\rangle.
\]
Then $L'\in\Pi_k^p=\Sigma_k^p$, so $\langle x,y_1\rangle\in L'\iff\exists z_1\forall z_2\cdots W$ accepts, for some polynomial-time $W$. Merging the two existential quantifiers $\exists y_1\exists z_1$ into one over the pair (padding to a common length), $L$ is defined by $k$ alternating quantifiers starting with $\exists$: $L\in\Sigma_k^p$. If $\mathsf P=\mathsf{NP}$, then $\mathsf{NP}=\mathsf{coNP}$ (as $\mathsf P$ is closed under complement), so $\mathsf{PH}=\Sigma_1^p=\mathsf{NP}=\mathsf P$.
\end{proof}

``The polynomial hierarchy does not collapse'' has become a standard working hypothesis, a little stronger than $\mathsf P\ne\mathsf{NP}$, and several theorems of the next chapters are of the form ``if this were true, the hierarchy would collapse''.

\section{Quantified formulas and games}\label{sec:pspace}

A \emph{quantified Boolean formula} is a sentence $Q_1x_1\,Q_2x_2\cdots Q_nx_n\ \psi$, where each $Q_i$ is $\exists$ or $\forall$, the $x_i$ range over $\{0,1\}$, and $\psi$ is a propositional formula in $x_1,\dots,x_n$. \textup{\textsc{Tqbf}} is the set of true ones. A language $B$ is \emph{$\mathsf{PSPACE}$-complete} if $B\in\mathsf{PSPACE}$ and $A\le_pB$ for every $A\in\mathsf{PSPACE}$.

\begin{theorem}[Stockmeyer--Meyer 1973]\label{thm:tqbf}
\textup{\textsc{Tqbf}} is $\mathsf{PSPACE}$-complete.
\end{theorem}

\begin{proof}
\emph{In $\mathsf{PSPACE}$.} To evaluate $Q_1x_1\cdots Q_nx_n\,\psi$, evaluate recursively the two sentences obtained by setting $x_1=0$ and $x_1=1$, and combine the answers by ``or'' or ``and''. The recursion has depth $n$; each level remembers one bit and where it is, and at the bottom $\psi$ is evaluated under a full assignment. The space is polynomial, since it is reused between the two branches.

\emph{Hardness.} Let $L\in\mathsf{PSPACE}$, decided by a one-tape machine $M$ in space $s=s(n)$, a polynomial. We may assume that $M$ uses only the cells $0,\dots,s-1$, and that before accepting it erases its tape and moves its head to cell $0$, so that there is a single accepting configuration $c_{\mathrm{acc}}$. A configuration is described by $O(s)$ Boolean variables, exactly as one row of the table of Theorem~\ref{thm:cooklevin}: variables $X_{i,a}$, $H_i$, $S_q$ for the cells $0\le i<s$. As there, with two vectors of variables $c$ and $c'$ we can write formulas of polynomial size in $s$:
\begin{itemize}
\item $\mathrm{valid}(c)$: the exactly-one clauses, saying that $c$ describes a configuration;
\item $\mathrm{step}(c,c')$: the move and frame clauses, with times $t$ and $t+1$ replaced by $c$ and $c'$, saying that $c'$ follows $c$ in one step;
\item $c=c'$: the conjunction of the equivalences of corresponding variables.
\end{itemize}
Let $\varphi_1(c,c')=\mathrm{valid}(c)\land\mathrm{valid}(c')\land\bigl(c=c'\lor\mathrm{step}(c,c')\bigr)$, and, with new vectors of variables $m$, $d$, $e$,
\[
\varphi_{2t}(c,c')=\exists m\ \forall d\ \forall e\ \Bigl[\bigl((d=c\land e=m)\lor(d=m\land e=c')\bigr)\to\varphi_t(d,e)\Bigr].
\]
We claim that $\varphi_t(c,c')$ holds if and only if $c$ and $c'$ are configurations and $c'$ is reachable from $c$ in at most $t$ steps. For $t=1$ this is the definition. For $2t$: if $c'$ is reachable from $c$ in at most $2t$ steps, take for $m$ a configuration in the middle of such a computation; for every $d,e$, the hypothesis of the implication says that $(d,e)$ is $(c,m)$ or $(m,c')$, and $\varphi_t(d,e)$ holds by induction. Conversely, if $\varphi_{2t}(c,c')$ holds with a witness $m$, take $(d,e)=(c,m)$ and then $(m,c')$: both $\varphi_t(c,m)$ and $\varphi_t(m,c')$ hold, and $c'$ is reachable in at most $t+t$ steps.

The trick is that $\varphi_t$ appears \emph{once} in $\varphi_{2t}$, the universal quantifier choosing which of the two halves to check; this is Savitch's recursion (Theorem~\ref{thm:savitch}) written as a formula. So $|\varphi_{2t}|=|\varphi_t|+O(s)$. The machine has at most $2^{cs}$ configurations, so if it accepts, it does so within $T=2^{cs}$ steps, and $\varphi_T$ has size $O(s^2)$ (the formula $\varphi_1$ has size $O(s^2)$, because of the exactly-one clauses for the head, and each of the $cs$ doublings adds $O(s)$). Substituting the constants of the initial configuration on $x$ for $c$ and of $c_{\mathrm{acc}}$ for $c'$, we get a sentence that is true if and only if $M$ accepts $x$. Finally, since $A\to\exists z\,B$ is equivalent to $\exists z\,(A\to B)$ when $z$ does not occur in $A$ (and similarly for $\forall$), all quantifiers can be moved to the front, giving a quantified Boolean formula, computed from $x$ in polynomial time.
\end{proof}

Without the universal quantifier, writing $\varphi_t(c,m)\land\varphi_t(m,c')$ would double the size at each step, and the formula would have exponential size. A single quantifier alternation saves an exponential: this is the flavour of the whole chapter.

To turn formulas into games we need them in a normal form. The next lemma, due to Grigori Tseitin (1968), costs only existential quantifiers.

\begin{lemma}\label{lem:tseitin}
For every propositional formula $\psi(x_1,\dots,x_n)$ there is a CNF formula $\chi$ in the variables $x_1,\dots,x_n$ and new variables $g_1,\dots,g_r$, with clauses of at most three literals, of size linear in $|\psi|$ and computed in polynomial time, such that for every assignment $a$ of the $x_i$: $\psi(a)$ is true if and only if $\exists g\ \chi(a,g)$ is true.
\end{lemma}

\begin{proof}
Write $\psi$ with $\neg$, $\land$, $\lor$, and introduce a variable $g$ for each subformula, $g=x_i$ for the atoms. For a subformula $g=\neg h$ add the clauses $(g\lor h)$, $(\neg g\lor\neg h)$; for $g=h\land h'$ the clauses $(\neg g\lor h)$, $(\neg g\lor h')$, $(g\lor\neg h\lor\neg h')$; for $g=h\lor h'$ the clauses $(g\lor\neg h)$, $(g\lor\neg h')$, $(\neg g\lor h\lor h')$; and finally the unit clause $(g_\psi)$ for the whole formula. The clauses of each subformula say exactly that $g$ equals the value of the subformula, given the values of $h$ and $h'$. So for each assignment $a$ the only $g$ satisfying all clauses but the last is the vector of values of the subformulas under $a$, and the last clause holds if and only if $\psi(a)$ is true.
\end{proof}

So in a quantified formula $Q_1x_1\cdots Q_nx_n\,\psi$ we may replace $\psi$ by $\exists g\,\chi$, with a 3-CNF matrix. Adding dummy variables, which appear in no clause, we may also make the quantifiers alternate, $\exists x_1\forall x_2\exists x_3\cdots\exists x_n$, starting and ending with $\exists$. In this form the sentence is a game: player E chooses $x_1$, player A chooses $x_2$, and so on; E wins if the chosen values satisfy the 3-CNF matrix. The sentence is true if and only if E has a winning strategy. Such \emph{formula games} are $\mathsf{PSPACE}$-complete by Theorem~\ref{thm:tqbf}, and they can be disguised as games on graphs.

In \emph{generalised geography}, a directed graph and a start vertex $b$ are given. A token stands on $b$. The two players, E first, alternately move the token along an arc to a vertex that has not been visited yet; a player who cannot move loses.

\begin{theorem}[Schaefer 1978]\label{thm:geography}
Deciding whether the first player has a winning strategy in generalised geography is $\mathsf{PSPACE}$-complete.
\end{theorem}

\begin{proof}
\emph{In $\mathsf{PSPACE}$.} A game lasts at most $|V|$ moves. A recursive search of the game tree (the current player wins if some move leads to a position where the other player loses) has depth at most $|V|$ and remembers one path; the space is polynomial.

\emph{Hardness.} Take a formula game $\exists x_1\forall x_2\cdots\exists x_n\,\chi$, with $n$ odd and $\chi$ a CNF with clauses $c_1,\dots,c_m$. Build the graph of Figure~\ref{fig:geography}: for each variable $x_i$ a diamond with a top $T_i$, two sides $L_i$ and $R_i$ and a bottom $B_i$, with arcs $T_i\to L_i$, $T_i\to R_i$, $L_i\to B_i$, $R_i\to B_i$; arcs $B_i\to T_{i+1}$ chaining the diamonds; a vertex $c$ with an arc $B_n\to c$; a vertex for each clause, with arcs $c\to c_j$; and, from the clause vertices, arcs $c_j\to R_i$ if $x_i$ is a literal of $c_j$, and $c_j\to L_i$ if $\neg x_i$ is. The start vertex is $T_1$.

The first $3n$ moves are almost forced. In diamond $i$, the move from $T_i$ to a side is a choice; the moves to $B_i$ and then to $T_{i+1}$ are forced. Each diamond takes three moves, so the choice at $T_i$ is made by E when $i$ is odd and by A when $i$ is even. Read the choice at $T_i$ as a value: $x_i$ is true if the token went through $L_i$. After the diamonds, at move $3n+1$, which is even since $n$ is odd, A moves from $c$ to a clause vertex $c_j$, of their choice. Then E must move to a side vertex $L_i$ or $R_i$ pointed to by $c_j$ that has not been visited, and from there A would have to move to $B_i$, which has been visited: A loses. If no such side vertex exists, E is stuck and loses.

Now the side vertex pointed to by a literal is unvisited exactly when the literal is true: for the literal $x_i$, the arc goes to $R_i$, which is unvisited when the token went through $L_i$, that is when $x_i$ is true; for $\neg x_i$, the arc goes to $L_i$, unvisited when $x_i$ is false. So after the diamonds E wins if and only if every clause has a true literal (A picks a clause with no true literal if there is one). Hence E wins the geography game if and only if E wins the formula game, and the graph has $4n+m+1$ vertices.
\end{proof}

\begin{figure}[ht]
\centering
\fig{ch10-geography}
\caption{Generalised geography built from the formula game $\exists x_1\forall x_2\exists x_3\ (x_1\lor\neg x_2\lor x_3)\land(\neg x_1\lor x_2)$. Going through $L_i$ means $x_i$ true. At the end, A picks a clause, and E must move to a side of a diamond that the token has not visited; the arcs of a clause point to the sides that are unvisited exactly when its literals are true.}
\label{fig:geography}
\end{figure}

Real games behave in the same way, with more engineering. Go, with a bound on the length of the game, is $\mathsf{PSPACE}$-complete; with the Japanese rule of ko, which allows very long games, John Robson proved in 1983 that it is $\mathsf{EXP}$-complete; and so is chess on $n\times n$ boards (all quoted). By the time hierarchy theorem, $\mathsf{EXP}$-complete problems are provably not in $\mathsf P$: here, for once, we know that no fast algorithm exists.

We can now keep the promise of Chapter~\ref{ch:automata}. For a nondeterministic finite automaton $N$ over an alphabet $\Sigma$, \emph{universality} asks whether $N$ accepts every word of $\Sigma^*$.

\begin{theorem}[Meyer--Stockmeyer 1972]\label{thm:nfauniversal}
Universality of nondeterministic finite automata is $\mathsf{PSPACE}$-complete.
\end{theorem}

\begin{proof}
\emph{In $\mathsf{PSPACE}$.} We show that non-universality is in $\mathsf{NPSPACE}$, which equals $\mathsf{PSPACE}$ by Savitch's theorem, and $\mathsf{PSPACE}$ is closed under complement. If $N$, with $k$ states, rejects some word, it rejects one of length less than $2^k$: along a shortest such word the sets of states reached after each prefix (the states of the subset construction, Theorem~\ref{thm:subset}) are distinct. A nondeterministic machine guesses the word letter by letter, keeping only the current set of states ($k$ bits) and a counter up to $2^k$ ($k$ bits), and accepts if it reaches a set without accepting states.

\emph{Hardness.} Let $L\in\mathsf{PSPACE}$. Its complement is in $\mathsf{PSPACE}$ too; let $M$ be a deterministic one-tape machine deciding $\bar L$ in space $s=s(n)$, using only the cells $0,\dots,s-1$, and staying still once it halts, as in Theorem~\ref{thm:cooklevin}. A configuration is written as a word of length $s+1$ over $\Gamma\cup Q$: the tape contents, with the state symbol inserted just before the scanned cell. An \emph{accepting history} of $M$ on $x$ is a word
\[
\#C_0\#C_1\#\cdots\#C_r\#,
\]
where $C_0$ is the initial configuration on $x$, each $C_{t+1}$ follows $C_t$, and $C_r$ contains $q_{\mathrm{acc}}$. In a history, the symbol at position $j$ of $C_{t+1}$ is determined by the symbols at positions $j-1,j,j+1,j+2$ of $C_t$ (reading $\#$ beyond the ends): if none of these four is a state symbol, it is the symbol at position $j$; otherwise the state symbol and the scanned symbol are among them and the transition decides it. This defines a finite table of \emph{legal} combinations of five symbols, depending only on $M$.

We build an NFA $N_x$ that accepts exactly the words that are \emph{not} accepting histories, as a union of NFAs, each looking for one kind of defect:
\begin{enumerate}
\item the word does not begin with $\#C_0\#$: a chain of $s+3$ states compares the beginning with the fixed word $\#C_0\#$, and accepts as soon as it finds a mismatch or the word ends;
\item the word contains no $q_{\mathrm{acc}}$, or does not end with $\#$: two states each;
\item some block between consecutive $\#$'s has a length different from $s+1$, or does not contain exactly one state symbol: the NFA guesses the block and counts, with $O(s)$ states;
\item some position violates the table: the NFA guesses a position $j$ in some block, remembers the four symbols at positions $j-1,\dots,j+2$, reads on, counting with $O(s)$ states, to the symbol at position $j$ of the next block ($s+2$ symbols after position $j$), and accepts if it is not legal for the four remembered ones.
\end{enumerate}
Each part has $O(s)$ states (times a constant depending on $M$) and is written in polynomial time. A word that has none of these defects is an accepting history; an accepting history has none of them. So $N_x$ is universal if and only if $M$ has no accepting history on $x$, that is, $x\notin\bar L$, that is, $x\in L$. This reduces $L$ to universality.
\end{proof}

The proof shows why NFAs are so much more concise than DFAs. The NFA has polynomially many states and checks a defect by \emph{guessing where it is}; a DFA would have to watch every position at once. Equivalence of NFAs, and of regular expressions, is $\mathsf{PSPACE}$-complete for the same reason.

\section{Counting}\label{sec:counting}

\begin{definition}\label{def:sharpp}
A function $f:\{0,1\}^*\to\N$ is in $\#\mathsf P$ if there are a polynomial $p$ and a polynomial-time machine $V$ with
\[
f(x)=\#\bigl\{y\in\{0,1\}^{p(|x|)}:V\text{ accepts }\langle x,y\rangle\bigr\}.
\]
A function $g$ is \emph{$\#\mathsf P$-hard} if every $f\in\#\mathsf P$ is computable in polynomial time by a machine allowed to ask the values of $g$ at unit cost; \emph{$\#\mathsf P$-complete} if moreover $g\in\#\mathsf P$.
\end{definition}

The counting version of an $\mathsf{NP}$ problem counts certificates: \#\textup{\textsc{Sat}} counts the satisfying assignments of a CNF formula, \#\textup{\textsc{Ham}} the Hamiltonian cycles of a graph, and so on. A counting oracle answers every $\mathsf{NP}$ question (is the count positive?), and a lot more.

\begin{proposition}\label{prop:sharpsat}
\#\textup{\textsc{Sat}} is $\#\mathsf P$-complete.
\end{proposition}

\begin{proof}
Let $f\in\#\mathsf P$, with $V$ and $p$. The formula $\varphi_x$ of the proof of Theorem~\ref{thm:cooklevin}, built from $V$, has exactly $f(x)$ satisfying assignments: a satisfying assignment determines the certificate written in the first row; conversely, once the certificate is fixed, all the variables of the table are determined, because $V$ is deterministic, so each accepted certificate gives exactly one satisfying assignment. So $f(x)$ is the number of models of $\varphi_x$, one question to the oracle.
\end{proof}

A reduction that preserves the number of solutions, like this one, is called \emph{parsimonious}. The interesting fact about $\#\mathsf P$ is that problems whose decision version is easy can be hard to count. The \emph{permanent} of an $n\times n$ matrix $A=(a_{ij})$ is
\[
\mathrm{perm}(A)=\sum_{\sigma\in S_n}\prod_{i=1}^na_{i\sigma(i)}.
\]
If $A$ is the $0$--$1$ matrix of a bipartite graph between $\{u_1,\dots,u_n\}$ and $\{v_1,\dots,v_n\}$ ($a_{ij}=1$ when $u_iv_j$ is an edge), each permutation contributes $1$ exactly when it is a perfect matching, so $\mathrm{perm}(A)$ counts perfect matchings. Deciding whether one exists is in $\mathsf P$ (Tome~C1, Chapter~6, and Tome~C2, Chapter~6).

\begin{theorem}[Valiant 1979]\label{thm:valiant}
Computing the permanent of $0$--$1$ matrices, that is, counting the perfect matchings of bipartite graphs, is $\#\mathsf P$-complete. \emph{(Quoted.)}
\end{theorem}

Valiant's proof reduces \#\textup{\textsc{Sat}} to the permanent of integer matrices with gadgets whose permanents cancel unwanted contributions, and then removes the negative entries by computing modulo large numbers. The determinant, which differs only by signs, is computed in $O(n^3)$ operations; signs make cancellations possible, and cancellations make Gaussian elimination work. Pólya's question of 1913 asked when signs can be attached so that the determinant equals the permanent; for bipartite graphs the answer is ``when the graph has a Pfaffian orientation'', characterised by Neil Robertson, Paul Seymour and Robin Thomas, in 1999, and independently by William McCuaig (quoted). Planar graphs have one, which is why the perfect matchings of planar graphs can be counted in polynomial time (Kasteleyn, and Temperley and Fisher, 1961, quoted), the source of the exact solution of the dimer model of statistical physics.

The best known exact algorithm for the permanent is still exponential, but much better than $n!$.

\begin{proposition}[Ryser 1963]\label{prop:ryser}
For an $n\times n$ matrix $A$,
\[
\mathrm{perm}(A)=\sum_{S\subseteq\{1,\dots,n\}}(-1)^{n-|S|}\prod_{i=1}^n\sum_{j\in S}a_{ij},
\]
which can be evaluated with $O(2^nn^2)$ arithmetic operations.
\end{proposition}

\begin{proof}
Expanding the product, $\prod_i\sum_{j\in S}a_{ij}=\sum_f\prod_ia_{if(i)}$, the sum over all maps $f:\{1,\dots,n\}\to S$. So the right-hand side is $\sum_f\bigl(\prod_ia_{if(i)}\bigr)\,\epsilon(f)$, the sum over all maps $f:\{1,\dots,n\}\to\{1,\dots,n\}$, with
\[
\epsilon(f)=\sum_{S\supseteq f(\{1,\dots,n\})}(-1)^{n-|S|}=\sum_{T\subseteq\{1,\dots,n\}\setminus I}(-1)^{n-|I|-|T|}=(1-1)^{n-|I|},
\]
where $I$ is the image of $f$ and $S=I\cup T$. This is $1$ if $I=\{1,\dots,n\}$, that is, if $f$ is a permutation, and $0$ otherwise (inclusion and exclusion, as in Tome~C1, Chapter~3). For each of the $2^n$ sets $S$, the product of $n$ sums takes $O(n^2)$ operations.
\end{proof}

Counting sits very high in the landscape. Toda's theorem says that the whole polynomial hierarchy is contained in $\mathsf P^{\#\mathsf P}$, the class of problems solvable in polynomial time with a counting oracle (Toda 1991, quoted): one counting question is worth any fixed number of quantifier alternations.

\section{The exponential time hypotheses}\label{sec:seth}

Brute force decides satisfiability of a formula with $n$ variables in time $2^n$ times a polynomial. For 3-SAT one can do better: the best known algorithms run in time about $1.31^n$, by refinements of an algorithm of Ramamohan Paturi, Pavel Pudlák, Michael Saks and Francis Zane (1998, quoted). For $k$-SAT the known algorithms run in time $2^{(1-c/k)n}$, approaching $2^n$ as $k$ grows. Russell Impagliazzo and Paturi proposed in 2001 to take this as a hypothesis.

\begin{definition}\label{def:eth}
The \emph{exponential time hypothesis} (ETH) is that there is $\delta>0$ such that 3-SAT cannot be decided in time $O(2^{\delta n})$, where $n$ is the number of variables. The \emph{strong exponential time hypothesis} (SETH) is that for every $\varepsilon>0$ there is $k$ such that $k$-SAT cannot be decided in time $O(2^{(1-\varepsilon)n})$.
\end{definition}

Both imply $\mathsf P\ne\mathsf{NP}$, and say much more: they fix the exponent. They are hypotheses, believed by many but not all; SETH in particular is much less certain than $\mathsf P\ne\mathsf{NP}$. Their value is in the reductions they allow. Under ETH, the sparsification lemma of Impagliazzo, Paturi and Zane (2001, quoted), which reduces a $k$-CNF formula to formulas with a linear number of clauses, shows that 3-colouring, independent set, Hamiltonian cycle and many others have no algorithms in time $2^{o(n)}$ (Exercise~\ref{ex:ethcolouring}). Under SETH, quadratic algorithms in $\mathsf P$ become optimal. The bridge is a simple problem: \textup{\textsc{Orthogonal vectors}} asks, given two sets $A$ and $B$ of $N$ vectors in $\{0,1\}^d$, whether some $a\in A$ and $b\in B$ have $a\cdot b=0$. It is solved in time $O(N^2d)$ by trying all pairs.

\begin{theorem}[Williams 2005]\label{thm:ov}
If for some $\varepsilon>0$ \textup{\textsc{Orthogonal vectors}} can be solved in time $O(N^{2-\varepsilon}\,d^{c})$ for a constant $c$, then SETH is false.
\end{theorem}

\begin{proof}
Let $\varphi$ be a $k$-CNF formula with $n$ variables (say $n$ even) and clauses $c_1,\dots,c_m$, where $m\le(2n)^k$. Split the variables into two halves $X_1$ and $X_2$ of $n/2$ variables. For each assignment $\alpha$ of $X_1$, let $a_\alpha\in\{0,1\}^m$ have $j$-th coordinate $0$ if some literal of $c_j$ on $X_1$ is made true by $\alpha$, and $1$ otherwise; define $b_\beta$ in the same way for the assignments $\beta$ of $X_2$ (Figure~\ref{fig:ov}). The assignment $(\alpha,\beta)$ satisfies $c_j$ if and only if $\alpha$ or $\beta$ makes one of its literals true, that is, if and only if $a_\alpha[j]=0$ or $b_\beta[j]=0$. So $(\alpha,\beta)$ satisfies $\varphi$ if and only if $a_\alpha\cdot b_\beta=0$. This gives two sets of $N=2^{n/2}$ vectors in dimension $d=m$, built in time $O(2^{n/2}nm)$, and an algorithm in time $O(N^{2-\varepsilon}d^c)$ decides $\varphi$ in time
\[
O\bigl(2^{(1-\varepsilon/2)n}(2n)^{kc}+2^{n/2}nm\bigr)=O\bigl(2^{(1-\varepsilon/4)n}\bigr)
\]
for every fixed $k$. So for every $k$, $k$-SAT is solved in time $O(2^{(1-\varepsilon/4)n})$, contradicting SETH for $\varepsilon/4$.
\end{proof}

\begin{figure}[ht]
\centering
\fig{ch10-ov}
\caption{The reduction of Theorem~\ref{thm:ov} on a formula with four variables and four clauses. Each assignment of half of the variables gives a vector with a $1$ (shaded) for each clause it does not satisfy; two halves combine into a satisfying assignment exactly when their vectors have no $1$ in common.}
\label{fig:ov}
\end{figure}

From orthogonal vectors, Arturs Backurs and Piotr Indyk built in 2015 two strings whose edit distance reveals whether an orthogonal pair exists, of length $N d^{O(1)}$ (quoted). So an algorithm for edit distance in time $O(n^{2-\varepsilon})$ would solve orthogonal vectors in time $O(N^{2-\varepsilon}d^{O(1)})$ and refute SETH. This is the theorem announced in Tome~C2, Chapter~5. Similar reductions explain why the longest common subsequence, the Fréchet distance between curves, and many other problems seem stuck at quadratic time. A whole field, \emph{fine-grained complexity}, classifies problems inside $\mathsf P$ in this way, by the exponents of their best algorithms and the hypotheses that would make them optimal.

\begin{beauty}
A single quantifier changes everything. ``There is a path'' is in $\mathsf{NL}$; ``there is a midpoint such that for both halves\dots'' compresses an exponentially long computation into a formula of quadratic size, and puts all of $\mathsf{PSPACE}$ inside a game. The polynomial hierarchy, games, the permanent and the fine-grained world are all about the same question: what is the price of a quantifier, of a count, of one more bit of exponent? Nobody can prove that any of these prices is positive; all the evidence says that they are, and the theory built on that evidence is coherent enough to predict, correctly so far, which problems will resist.
\end{beauty}

\section{Exercises}

\begin{exercise}[$\star$]\label{ex:minformula}
Show that the language of pairs $(\varphi,k)$, where $\varphi$ is a propositional formula and some formula of size at most $k$ is equivalent to $\varphi$, is in $\Sigma_2^p$. Show that its complement is in $\Pi_2^p$, and explain why it is not obviously in $\mathsf{NP}$ or in $\mathsf{coNP}$.
\end{exercise}

\begin{exercise}[$\star$]\label{ex:phpspace}
Show that if $\mathsf{PH}=\mathsf{PSPACE}$, then the polynomial hierarchy collapses to some level $\Sigma_k^p$.
\end{exercise}

\begin{exercise}[$\star$]\label{ex:qbf2cnf}
Show that \textup{\textsc{Tqbf}} restricted to sentences whose matrix is a 2-CNF formula can be decided in polynomial time. (You may quote that 2-SAT is in $\mathsf P$, and handle the quantifiers one at a time from the inside.)
\end{exercise}

\begin{exercise}[$\star\star\star$]\label{ex:undirectedgeo}
In \emph{undirected vertex geography}, the graph is undirected. Show that the first player wins from $b$ if and only if every maximum matching of the graph covers $b$, and deduce that this game can be decided in polynomial time.
\end{exercise}

\begin{exercise}[$\star$]\label{ex:permallones}
Show that the permanent of the $n\times n$ matrix with zeros on the diagonal and ones elsewhere is the number $D_n$ of derangements (permutations without fixed points). Deduce from Proposition~\ref{prop:ryser} a formula for $D_n$ as a sum of $n+1$ terms, and check it for $n\le5$ against $D_n=n!\sum_{i=0}^n(-1)^i/i!$.
\end{exercise}

\begin{exercise}[$\star$]\label{ex:permmod2}
Show that the parity of the number of perfect matchings of a bipartite graph can be computed in polynomial time.
\end{exercise}

\begin{exercise}[$\star$]\label{ex:sharpdnf}
Satisfiability of formulas in disjunctive normal form is in $\mathsf P$. Show that counting their satisfying assignments is $\#\mathsf P$-complete.
\end{exercise}

\begin{exercise}[$\star\star$]\label{ex:parsimonious}
Show that the reduction of Theorem~\ref{thm:threesat} is not parsimonious, and modify it so that it is.
\end{exercise}

\begin{exercise}[$\star$]\label{ex:ethcolouring}
Assume ETH, and the sparsification lemma in the form: if ETH holds, there is $\delta'>0$ such that 3-SAT on formulas with $n$ variables and at most $Cn$ clauses, for some constant $C$, cannot be decided in time $O(2^{\delta'n})$. Show that $3$-colourability of graphs with $N$ vertices cannot be decided in time $2^{o(N)}$.
\end{exercise}

\begin{exercise}[$\star$]\label{ex:ovsmall}
Show that \textup{\textsc{Orthogonal vectors}} can be solved in time $O(2^dd+Nd)$. Why does this not contradict Theorem~\ref{thm:ov}?
\end{exercise}

\section{Hints}

\begin{hint}{ex:minformula}
Guess $\psi$; for all assignments compare. For the complement, swap the quantifiers.
\end{hint}
\begin{hint}{ex:phpspace}
\textup{\textsc{Tqbf}} would lie in some $\Sigma_k^p$.
\end{hint}
\begin{hint}{ex:qbf2cnf}
Eliminating an innermost $\exists x$ from a 2-CNF gives a 2-CNF (resolve on $x$); eliminating an innermost $\forall x$ amounts to deleting $x$ from all clauses.
\end{hint}
\begin{hint}{ex:undirectedgeo}
If a maximum matching $M$ covers $b$, the first player always moves along $M$. If some maximum matching misses $b$, the second player answers along it.
\end{hint}
\begin{hint}{ex:permallones}
A permutation contributes $1$ exactly when it has no fixed point.
\end{hint}
\begin{hint}{ex:permmod2}
Modulo $2$, signs do not matter.
\end{hint}
\begin{hint}{ex:sharpdnf}
The negation of a CNF formula is a DNF formula with the same number of variables.
\end{hint}
\begin{hint}{ex:parsimonious}
Make each new variable equal to a disjunction of literals.
\end{hint}
\begin{hint}{ex:ethcolouring}
Count the vertices of the graph of Theorem~\ref{thm:threecol}.
\end{hint}
\begin{hint}{ex:ovsmall}
Tabulate which subsets of coordinates are the supports of vectors of $B$.
\end{hint}

\section{Solutions}

\begin{solution}{ex:minformula}
\idea One existential guess for the small formula, one universal check over all assignments.
\writeup $(\varphi,k)$ is in the language if and only if there is a formula $\psi$ of size at most $k$ (a word of length polynomial in $|\varphi|+k$; we may assume $k\le|\varphi|$, since $\varphi$ itself is a candidate) such that for every assignment $a$ of the variables, $\varphi(a)=\psi(a)$; the check takes polynomial time. That is a $\Sigma_2^p$ definition. The complement reads: for every $\psi$ of size at most $k$ there is an assignment $a$ with $\varphi(a)\ne\psi(a)$, a $\Pi_2^p$ definition. A certificate for membership would have to convince us both of the existence of $\psi$ and of an equivalence, which is a $\mathsf{coNP}$ statement; nobody knows short certificates for both at once.
\end{solution}

\begin{solution}{ex:phpspace}
\idea A complete problem pins the class to one level.
\writeup If $\mathsf{PH}=\mathsf{PSPACE}$, then \textup{\textsc{Tqbf}} is in $\mathsf{PH}$, so it is in $\Sigma_k^p$ for some $k$. Every $L\in\mathsf{PSPACE}$ reduces to \textup{\textsc{Tqbf}} by some $f$, and $x\in L$ if and only if $f(x)$ is in \textup{\textsc{Tqbf}}, which is a $\Sigma_k^p$ condition (compute $f(x)$ inside the verifier). So $\mathsf{PSPACE}\subseteq\Sigma_k^p$, and $\mathsf{PH}=\Sigma_k^p$.
\end{solution}

\begin{solution}{ex:qbf2cnf}
\idea Quantifier elimination preserves 2-CNF, and each elimination is a syntactic operation.
\writeup Process the quantifiers from the innermost. If the innermost is $\exists x$ and the matrix is a 2-CNF $\chi$, then $\exists x\,\chi$ is equivalent to the conjunction of the clauses of $\chi$ not containing $x$, together with all resolvents $a\lor b$ of a clause $x\lor a$ with a clause $\neg x\lor b$ (Tome~C1, Chapter~1): if all of these hold, one of the two values of $x$ satisfies the clauses containing it. The result is again a 2-CNF (clauses of two literals, possibly fewer), with at most quadratically many clauses in the number of variables, so its size stays polynomial. If the innermost is $\forall x$, then $\forall x\,\chi$ holds if and only if $\chi$ holds for $x=0$ and for $x=1$; a clause $x\lor a$ must then hold with $x=0$, so it becomes $a$, and similarly for $\neg x\lor a$: delete $x$ and $\neg x$ from all clauses (an empty clause means false). After all quantifiers are eliminated, a 2-CNF without variables remains, true if and only if it has no empty clause. Each step takes polynomial time, and clauses never exceed two literals, so the number of distinct clauses stays $O(n^2)$.
\end{solution}

\begin{solution}{ex:undirectedgeo}
\idea Matchings pair each of the opponent's moves with an answer.
\writeup Let $G$ be the graph and $b$ the start. Suppose every maximum matching covers $b$; fix one, $M$. The first player always moves along $M$: from the current vertex $v$ to its partner. We claim this is always possible. At the start, $b$ is covered. Suppose the first player has so far moved along $M$, the visited vertices are $b=v_0,v_1,\dots,v_{2r}$ with $v_{2i}v_{2i+1}\in M$ for $i<r$, and the second player has just moved to $v_{2r}$. If $v_{2r}$ were not covered by $M$, then $v_0v_1\dots v_{2r}$ would be an alternating path from $b$, starting with an edge of $M$ and ending at an uncovered vertex; exchanging the edges along it gives a maximum matching (same size) not covering $b$, a contradiction. And the partner of $v_{2r}$ is not visited, since the visited vertices are paired among themselves except $v_{2r}$. So the first player can always move, and the second player eventually gets stuck.

Conversely, suppose some maximum matching $M$ does not cover $b$. The second player answers each move along $M$. The first move goes to some $v_1$; $v_1$ is covered by $M$, otherwise adding the edge $bv_1$ would give a larger matching. In general, after moves $b,v_1,\dots,v_{2r+1}$ with $v_{2i+1}v_{2i+2}\in M$, if $v_{2r+1}$ were uncovered, the path $b,v_1,\dots,v_{2r+1}$ would be augmenting (both ends uncovered, edges alternating), contradicting maximality (Tome~C1, Chapter~6). So the second player always has an answer, and wins. The criterion is checked in polynomial time: compute the size $\nu(G)$ of a maximum matching and the size $\nu(G-b)$; every maximum matching covers $b$ if and only if $\nu(G-b)<\nu(G)$.
\end{solution}

\begin{solution}{ex:permallones}
\idea Permanents of $0$--$1$ matrices count permutations avoiding the zeros.
\writeup A permutation $\sigma$ contributes $\prod_ia_{i\sigma(i)}$, which is $1$ when $\sigma(i)\ne i$ for every $i$ and $0$ otherwise, so the permanent is $D_n$. In Ryser's formula, for $|S|=s$, row $i$ has sum $s-1$ if $i\in S$ and $s$ otherwise, so the product is $(s-1)^ss^{n-s}$ and
\[
D_n=\sum_{s=0}^n(-1)^{n-s}\binom ns(s-1)^ss^{n-s}.
\]
(with $0^0=1$). For $n=1,\dots,5$ this gives $0,1,2,9,44$, which agrees with $n!\sum_{i=0}^n(-1)^i/i!$. For instance for $n=3$ the four terms are $0$, $0$, $-6$ and $8$, giving $2$: the two cyclic permutations of three elements.
\end{solution}

\begin{solution}{ex:permmod2}
\idea Signs disappear modulo $2$, so the permanent is a determinant there.
\writeup Since $-1\equiv1\pmod2$, $\mathrm{perm}(A)\equiv\det(A)\pmod2$. The determinant of a $0$--$1$ matrix modulo $2$ is computed by Gaussian elimination over the field $\mathbb F_2$ in $O(n^3)$ operations. So the parity of the number of perfect matchings is computable in polynomial time, while the number itself is $\#\mathsf P$-complete; Valiant's theorem is about the integers, and modulo $2$ it fails. (Modulo $2^k$ the permanent is still polynomial, by a theorem of Valiant; modulo an odd prime it is hard.)
\end{solution}

\begin{solution}{ex:sharpdnf}
\idea Counting models of a DNF is counting non-models of a CNF.
\writeup Satisfiability of a DNF is easy: it is satisfiable if and only if some term contains no complementary pair of literals. Counting is in $\#\mathsf P$. For hardness, given a CNF $\varphi$ with $n$ variables, its negation is, by De Morgan's laws, a DNF $\psi$ of the same size, and the number of models of $\varphi$ is $2^n$ minus the number of models of $\psi$. One oracle question to \#\textup{\textsc{Dnf}} computes \#\textup{\textsc{Sat}}, which is $\#\mathsf P$-complete.
\end{solution}

\begin{solution}{ex:parsimonious}
\idea Free choices of auxiliary variables multiply the count; define them by equivalences instead.
\writeup In Theorem~\ref{thm:threesat}, for the clause $\ell_1\lor\ell_2\lor\ell_3\lor\ell_4$, the new clauses are $(\ell_1\lor\ell_2\lor z)\land(\neg z\lor\ell_3\lor\ell_4)$. If all four literals are true, both values of $z$ satisfy them: one model becomes two. To make the reduction parsimonious, force $z_1\leftrightarrow(\ell_1\lor\ell_2)$ and $z_{i}\leftrightarrow(z_{i-1}\lor\ell_{i+1})$, and require $z_{k-2}\lor\ell_k$. An equivalence $z\leftrightarrow(u\lor v)$ is the 3-CNF $(\neg z\lor u\lor v)\land(z\lor\neg u)\land(z\lor\neg v)$. For each assignment of the original variables the $z_i$ are determined, and the last clause holds exactly when the original clause does. So each model of the original formula extends in exactly one way.
\end{solution}

\begin{solution}{ex:ethcolouring}
\idea A linear-size reduction transports a linear exponent.
\writeup Given a 3-CNF formula with $n$ variables and $m\le Cn$ clauses, the graph of Theorem~\ref{thm:threecol} has $N=3+2n+6m\le(2+6C)n+3$ vertices. An algorithm for $3$-colourability in time $2^{o(N)}$ would decide such formulas in time $2^{o(n)}$, plus the polynomial time of the reduction, which is less than $2^{\delta'n}$ for large $n$. This contradicts the sparsified form of ETH. Without sparsification, $m$ could be of order $n^3$ and the argument would fail.
\end{solution}

\begin{solution}{ex:ovsmall}
\idea When the dimension is small, the number of possible vectors, $2^d$, is smaller than the number of pairs.
\writeup Mark in a table of size $2^d$ the vectors of $B$ (as subsets of $\{1,\dots,d\}$). Compute, for every subset $T$, whether some marked set is contained in $T$: by dynamic programming over the coordinates, in time $O(2^dd)$: after processing coordinates $1,\dots,i$, the entry of $T$ says whether some marked set agrees with $T$ outside $\{1,\dots,i\}$ and is contained in $T$ on $\{1,\dots,i\}$. Then $a\in A$ is orthogonal to some $b\in B$ if and only if some marked set is contained in the complement of the support of $a$: one lookup per vector of $A$. In the reduction of Theorem~\ref{thm:ov}, $d=m$ is at least $n$ while $N=2^{n/2}$, so $2^d\ge N^2$: the dimension is too large for this algorithm to help. Orthogonal vectors is believed hard when $d$ is somewhat larger than $\log N$, for example $d=(\log N)^2$.
\end{solution}

\section*{Sources and further reading}

\begin{itemize}
\item \emph{A research paper to read.} L.~G.~Valiant, \emph{The complexity of computing the permanent}, Theoretical Computer Science 8 (1979), 189--201. The paper that founded counting complexity. It begins by defining $\#\mathsf P$ and comparing the permanent with the determinant, and points out what was new: a problem whose decision version is easy and whose counting version is as hard as any counting problem. The main reduction goes from counting the models of a 3-CNF formula to the permanent of a matrix with small integer entries, some negative, viewed as a weighted directed graph whose cycle covers are counted with weights; it is built from gadgets for the variables, for the clauses, and for the connections between them, small matrices whose minors have prescribed permanents. Checking the properties of these gadgets by hand, with Proposition~\ref{prop:ryser} or by expansion, is a good exercise. The negative and large entries are then removed, by computing modulo a large number and by replacing weighted edges with small unweighted graphs, to reach $0$--$1$ matrices. Read it with the companion paper, L.~G.~Valiant, \emph{The complexity of enumeration and reliability problems}, SIAM Journal on Computing 8 (1979), 410--421, which shows that many other counting problems are complete, among them counting the models of monotone 2-CNF formulas, whose satisfiability is trivial.
\item \emph{The stories.} A.~R.~Meyer and L.~J.~Stockmeyer, \emph{The equivalence problem for regular expressions with squaring requires exponential space}, 13th Annual Symposium on Switching and Automata Theory (1972), 125--129; L.~J.~Stockmeyer and A.~R.~Meyer, \emph{Word problems requiring exponential time}, Proceedings of the Fifth ACM Symposium on Theory of Computing (1973), 1--9; L.~J.~Stockmeyer, \emph{The polynomial-time hierarchy}, Theoretical Computer Science 3 (1976), 1--22; G.~Pólya, Aufgabe 424, Archiv der Mathematik und Physik (3) 20 (1913), 271.
\item T.~J.~Schaefer, \emph{On the complexity of some two-person perfect-information games}, Journal of Computer and System Sciences 16 (1978), 185--225; D.~Lichtenstein and M.~Sipser, \emph{Go is polynomial-space hard}, Journal of the ACM 27 (1980), 393--401; A.~S.~Fraenkel and D.~Lichtenstein, \emph{Computing a perfect strategy for $n\times n$ chess requires time exponential in $n$}, Journal of Combinatorial Theory, Series A 31 (1981), 199--214; J.~M.~Robson, \emph{The complexity of Go}, Proceedings of the IFIP 9th World Computer Congress (1983), 413--417; A.~S.~Fraenkel, E.~R.~Scheinerman and D.~Ullman, \emph{Undirected edge geography}, Theoretical Computer Science 112 (1993), 371--381.
\item G.~S.~Tseitin, \emph{On the complexity of derivation in propositional calculus}, Studies in Constructive Mathematics and Mathematical Logic, Part~2 (1968), 115--125; C.~Umans, \emph{The minimum equivalent DNF problem and shortest implicants}, Journal of Computer and System Sciences 63 (2001), 597--611; S.~Toda, \emph{PP is as hard as the polynomial-time hierarchy}, SIAM Journal on Computing 20 (1991), 865--877; H.~J.~Ryser, \emph{Combinatorial Mathematics}, Mathematical Association of America, 1963; N.~Robertson, P.~D.~Seymour and R.~Thomas, \emph{Permanents, Pfaffian orientations, and even directed circuits}, Annals of Mathematics 150 (1999), 929--975.
\item R.~Impagliazzo and R.~Paturi, \emph{On the complexity of $k$-SAT}, Journal of Computer and System Sciences 62 (2001), 367--375; R.~Impagliazzo, R.~Paturi and F.~Zane, \emph{Which problems have strongly exponential complexity?}, Journal of Computer and System Sciences 63 (2001), 512--530; R.~Williams, \emph{A new algorithm for optimal 2-constraint satisfaction and its implications}, Theoretical Computer Science 348 (2005), 357--365; A.~Backurs and P.~Indyk, \emph{Edit distance cannot be computed in strongly subquadratic time (unless SETH is false)}, SIAM Journal on Computing 47 (2018), 1087--1097; V.~Vassilevska Williams, \emph{On some fine-grained questions in algorithms and complexity}, Proceedings of the International Congress of Mathematicians (2018), 3447--3487.
\item S.~Arora and B.~Barak, \emph{Computational Complexity: A Modern Approach}, Cambridge University Press, 2009, Chapters~4, 5, 9 and~17; M.~Sipser, \emph{Introduction to the Theory of Computation}, 3rd edition, Cengage, 2013, Chapter~8.
\end{itemize}
